    \chapter{Applications of Differential Calculus}
    \section{Applications to Geometry}
    \begin{definitionenv}
        Suppose $f$ is a scalar field on $S\subseteq \R^n$ and $c\in \R$ is a constant. We call the set $L(c) := \{\bm x \in S: f(\bm x) = c\}$ a \textbf{level set} of $f$. In particular, if $n=2$, we call $L(c)$ a \textbf{level curve} of $f$; if $n=3$, we call $L(c)$ a \textbf{level surface} of $f$.
    \end{definitionenv}
    \begin{exampleenv}
        Suppose $f:\R^2 \to \R$. For any curve $\bm r(t) $ in $\R^2$, if we let $g$ be defined as $g(t) = f[\bm r(t)]$, then $$g'(t) = \nabla f(\bm r(t)) \cdot \bm r'(t) = \| \bm r'(t)\| \left( \nabla f(\bm r(t)) \cdot T(t) \right).$$
        Particularly, if this curve happens to be the level curve given by $f(x,y)  = c$. Suppose level curve $L(c) := \{ (x,y) : f(x,y) = c\}$ is also described by $\bm r_c(t)$. Then $g(t) = f[\bm r_c(t)] = c$ for any $t$, which means $g'(t) = 0$ for any $t$. Thus
        $$
        \| \bm r_c'(t)\| \left( \nabla f(\bm r_c(t)) \cdot T(t) \right) = 0.
        $$
        On this level curve, $\nabla f \cdot T = 0$, and we say: the gradient vector is always normal to the curve $C$; the directional derivative of $f$ along any level curve of $f$ is zero. By Cauchy-Schwarz inequality, we have
        $$ |\nabla f(x,y) \cdot \bm y | \le \|\nabla f(x,y)\| \|\bm y\|, $$
        maximized when $\bm y$ is parallel to $\nabla f(x,y)$, i.e., $\bm y$ is perpendicular to the tangent vector of a level curve.
    \end{exampleenv}
    \begin{exampleenv}
        Suppose $f:\R^3 \to \R$ differentiable. Suppose $L(c)$ is a level surface and $\bm a$ is a point on the level surface. Let $\bm r (t)$ describes a curve that passes through $\bm a$ and lies completely with in $L(c)$. If we consider all curves passing through $\bm a $ and living within $L(c)$, then the tangent vectors of all therse curves at $\bm a$ determin/live in a plane that passes throught $\bm a$ and has a normal vector equal to $\nabla f(\bm a)$ provided that $\nabla f(\bm a) \ne \bm 0$. This plane is can be described by the equation
        $$ \nabla f(\bm a) \cdot (\bm x - \bm a) = 0. $$
    \end{exampleenv}

    \section{Applications to Optimization Problems}
    \begin{definitionenv}
        A scalar field is said to
        \begin{itemize}
            \item have an \textbf{absolute (or global maximum)} at $\bm a \in S \subseteq \R^n$ if $f(\bm a) \ge f(\bm x)$ for any $\bm x \in S$, in which case we also call $f(\bm a)$ the \textbf{absolute maximum value} of $f$ on $S$;
            \item have an \textbf{absolute (or global minimum)} at $\bm a \in S$ if $f(\bm a) \le f(\bm x)$ for any $\bm x \in S$, in which case we also call $f(\bm a)$ the \textbf{absolute minimum value} of $f$ on $S$;
            \item have a \textbf{relative (or local maximum)} at $\bm a \in S$ if there exists $r>0$ such that $f(\bm a) \ge f(\bm x)$ for any $\bm x \in B(\bm a; r)\subseteq S$, in which case we also call $f(\bm a)$ the \textbf{local maximum value} of $f$ at $\bm a$;
            \item have a \textbf{relative (or local minimum)} at $\bm a \in S$ if there exists $r>0$ such that $f(\bm a) \le f(\bm x)$ for any $\bm x \in B(\bm a; r)\subseteq S$, in which case we also call $f(\bm a)$ the \textbf{local minimum value} of $f$ at $\bm a$.
        \end{itemize}
        If $f$ has either a relative maximum or a relative minimum at $\bm a$, then we say $f$ has an \textbf{extremum} at $\bm a$.
    \end{definitionenv}
    \begin{theoremenv}
        If $f$ has an extremum at $\bm a\in \operatorname{int} S$ and $f$ is differentiable at $\bm a$, then $\nabla f(\bm a) = \bm 0$.
    \end{theoremenv}
    \begin{remarkenv}
        If $f$ is differentiable at $\bm a$ and $\nabla f(\bm a) = \bm 0$, then $f$ does not necessarily have an extremum at $\bm a$. Assume $z=f(x,y) = x^2 - y^2$, then $\nabla f(0,0) = \bm 0$ but $f$ does not have an extremum at $(0,0)$ since $f$ can be positive or negative in any neighborhood of $(0,0)$.
    \end{remarkenv}
    \begin{definitionenv}
        Suppose scalar field $f$ is differentiable at $\bm a$. If $\nabla f(\bm a) = \bm 0$, then we say $\bm a$ is a \textbf{stationary point} of $f$. A stationary point is called a \textbf{saddle point} if every open $n$-ball $B(\bm a)$ contains both points $\bm x$ with $f(\bm x) > f(\bm a)$ and points $\bm y$ with $f(\bm y) < f(\bm a)$.
    \end{definitionenv}
    \begin{remarkenv}
        Suppose $f$ has a stationary point at $\bm a$, then the nature of this stationary point is determined by the sign of $f(\bm x) - f(\bm a)$ for $\bm x$ close to $\bm a$. Suppose $\bm x = \bm a + \bm y$, then
        $$ f(\bm a + \bm y) - f(\bm a) = \nabla f(\bm a) \cdot \bm y + \| \bm y\| E(\bm a, \bm y),$$
        where $E(\bm a, \bm y) \to 0$ as $\bm y \to \bm 0$.
    \end{remarkenv}
    \begin{theoremenv}
        Suppose $f$ is a scalar field with continuous second-order partial derivatives $D_{ij} f$ in an open $n$-ball $B(\bm a)$. Then for all $\bm y \in \R^n$ such that $(\bm a + \bm y) \in B(\bm a)$, we have
        $$ 
        f(\bm a + \bm y) - f(\bm a ) = \nabla f(\bm a) \cdot \bm y + \frac{1}{2!} \bm y H(\bm a + c \bm y) \bm y^T,
        $$
        where $c\in (0,1)$ and $H(\bm x)$ is called the \textbf{Hessian matrix} whose $ij$-th entry is given by $D_{ij} f(\bm x) = \displaystyle \frac{\partial^2 f}{\partial x_i \partial x_j}(\bm x)$. Note that here $\bm y$ is treated as a row vector. We can also write the second-order Taylor expansion as
        $$
        f(\bm a + \bm y) = f(\bm a) + \nabla f(\bm a) \cdot \bm y + \frac{1}{2!} \bm y H(\bm a) \bm y^T + \| \bm y \|^2 E_2(\bm a, \bm y),
        $$
        where $E_2(\bm a, \bm y) \to \bm 0$ as $\bm y \to \bm 0$.
    \end{theoremenv}
    \begin{exampleenv}
        Consider the function
        $$
        f(x,y) = \begin{cases}
            xy\dfrac{x^2-y^2}{x^2+y^2} & (x,y) \ne (0,0) \\[6pt]
            0 & (x,y) = (0,0)
        \end{cases}.
        $$
        By definition,
        $$ \begin{aligned}
            D_1 f(0,y) &= \lim_{h\to 0} \frac{f(h,y) - f(0,y)}{h} = \lim_{h\to 0} y \frac{h^2 - y^2}{h^2 + y^2} = -y, \\ D_{2,1}f(0,0) &= D_2(D_1f)(0,0) = \lim_{k \to 0} \frac{D_1 f(0,k) - D_1 f(0,0)}{k} = -1.
        \end{aligned}$$
        Analogously, we have
        $$
        D_2 f(x,0) = x, \quad D_{1,2} f(0,0) = 1.
        $$
        Though $D_{1,2} f(0,0)$ and $D_{2,1} f(0,0)$ both exist, they are not equal since we do not have the continuity of $D_{1,2} f$ and $D_{2,1} f$ at $(0,0)$: note
        $$
        D_{2,1}f(x,y)
        =
        \frac{x^6+9x^4y^2-9x^2y^4-y^6}{(x^2+y^2)^3},
        \qquad (x,y)\neq(0,0),
        $$
        and we have
        $$ D_{2,1}f(0,0)=-1 \ne 1 =\frac{x^6}{x^6} =\lim_{x \to 0} D_{2,1}f(x,0)\quad \text{along } y=0. $$
        This example shows that the continuity of second-order partial derivatives is necessary for the equality of mixed partial derivatives and motivates the following theorem.
    \end{exampleenv}
    \begin{theoremenv}\label{thm: mixed partials}
        Suppose \(f\) is a scalar field on an open set \(S\subseteq \R^n\), and suppose
        the first-order partial derivatives \(D_i f\) and \(D_j f\) exist on \(S\).
        Assume further that the second-order partial derivatives \(D_{ij}f\) and
        \(D_{ji}f\) exist on \(S\) and are continuous at \(\bm a\in S\). Then
        \[
            D_{ij}f(\bm a)=D_{ji}f(\bm a).
        \]
    \end{theoremenv}
    \begin{remarkenv}
        One can actually prove a \textbf{stronger} version of the above theorem, which only requires the continuity of \textbf{one of the two} second order partial derivatives \(D_{i,j}f\) and \(D_{j,i}f\) at \(\bm a\).
    \end{remarkenv}
    \begin{lemmaenv}
        By Theorem~\ref{thm: mixed partials}, we know the symmetric Hessian matrix has an orthonormal eigenbasis and $\bm y H(\bm a + c \bm y) \bm y^T$ is a quadratic form associated with a unique symmetric matrix. If we know all the eigenvalues of $H(\bm a)$ are positive, then
        $$ f(\bm a + \bm y) - f(\bm a) = \frac{1}{2} \bm y H(\bm a) \bm y^T + \| \bm y\|^2 E_2(\bm a, \bm y) > 0 $$
        for $\bm y$ close enough to $\bm 0$.
    \end{lemmaenv}
    \begin{proofenv}
        Fix $\bm y$ and define $g(u) = f(\bm a + u\bm y)$ for $u \in [-1,1]$. Then
        $$ f(\bm a + \bm y ) = g(1) - g(0) = g'(0) + \frac{1}{2} g''(c)$$
        where $c\in (0,1)$ by the second-order Taylor formula for $g$ on $[0,1]$ with the Lagrange form of the remainder.

        Let $\bm r = \bm a + u \bm y$, then $\bm r'(u) = \bm y$ and
        $$ g'(u) = \nabla f(\bm r) \cdot \bm r'(u) = \sum_{j=1} D_j f(\bm r) y_j. $$
        Thus
        $$ g'(0) = \sum_{j=1}^{n=1} D_j f(\bm a) y_j = \nabla f(\bm a) \cdot \bm y = 0. $$

    \end{proofenv}
    % \begin{corollaryenv}
    %     Suppose $f$ is a scalar field on $\R^2$ with $D_1 f, D_2 f, D_{12} f, D_{21} f$ exist on an open set $S$. If $(a,b) \in S$ and $D_{12} f$ and $D_{21} f$ are both continuous at $(a,b)$, then $D_{12} f(a,b) = D_{21} f(a,b)$.
    % \end{corollaryenv}
    \begin{lemmaenv}
        Suppose $A\in \mathbf M_{n\times n}(\R)$ is a symmetric matrix and $\bm y $ is a row vector in $\R^n$. Let $Q(\bm y) = \bm y A \bm y^T$. Then
        \begin{enumerate}
            \item $Q(\bm y) > 0$ for any $\bm y \ne \bm 0$ if and only if all the eigenvalues of $A$ are positive;
            \item $Q(\bm y) < 0$ for any $\bm y \ne \bm 0$ if and only if all the eigenvalues of $A$ are negative;
            \item $Q(\bm y)$ takes both positive and negative values for $\bm y \ne \bm 0$ if and only if $A$ has both positive and negative eigenvalues.
        \end{enumerate}
    \end{lemmaenv}
    \begin{theoremenv}
        Suppose $f$ is a scalar field with continuous second-order partial derivatives $D_{ij} f$ in an open $n$-ball $B(\bm a)$ and Let $H(\bm a)$ denote the Hessian matrix of $f$ at $\bm a$, which is a stationary point of $f$. Then
        \begin{enumerate}
            \item If all the eigenvalues of $H(\bm a)$ are positive, i.e., $H(\bm a)$ is positive definite, then $f$ has a local minimum at $\bm a$;
            \item If all the eigenvalues of $H(\bm a)$ are negative, i.e., $H(\bm a)$ is negative definite, then $f$ has a local maximum at $\bm a$;
            \item If $H(\bm a)$ has both positive and negative eigenvalues, then $f$ has a saddle point at $\bm a$.
        \end{enumerate}
    \end{theoremenv}
    \begin{corollaryenv}
        Let \(\bm a\) be a stationary point of a scalar field \(f(x_1,x_2)\) with continuous second order partial derivatives in a 2-ball \(B(\bm a)\). Let
        \[
        A=D_{1,1}f(\bm a),\qquad B=D_{1,2}f(\bm a),\qquad C=D_{2,2}f(\bm a),
        \]
        and let
        \[
        \Delta=\det H(\bm a)
        =\det
        \begin{pmatrix}
        A & B\\
        B & C
        \end{pmatrix}
        =AC-B^2.
        \]
        Then we have:
        \begin{enumerate}[label=(\alph*)]
            \item If $\Delta > 0$ and $A > 0$, then $f$ has a local minimum at $\bm a$;
            \item If $\Delta > 0$ and $A < 0$, then $f$ has a local maximum at $\bm a$;
            \item If $\Delta < 0$, then $f$ has a saddle point at $\bm a$;
            \item If $\Delta = 0$, then the test is inconclusive.
        \end{enumerate}
    \end{corollaryenv}
    \begin{remarkenv}
        The determinant test is only applicable for $n=2$. For $n\ge 3$, we need to check the signs of all the eigenvalues of the Hessian matrix to determine the nature of a stationary point.
    \end{remarkenv}
    \begin{theoremenv}
        Given $\bm a = (a_1, \ldots, a_n), \bm b = (b_1, \ldots, b_n) \in \R^n$ with $a_i < b_i$ for $i=1, \ldots, n$. We define $\left[ \bm a, \bm b \right] := [a_1, b_1] \times \cdots \times [a_n, b_n]$. If a scalar field $f$ is continuous at every $\bm x \in [\bm a, \bm b]$, then
        \begin{itemize}
            \item $f$ is bounded on $[\bm a, \bm b]$, i.e., there exists $M>0$ such that $|f(\bm x)| \le M$ for any $\bm x \in [\bm a, \bm b]$;
            \item There exists $\bm c$ and $\bm d$ such that
            $$ f(\bm c) = \sup_{\bm x \in [\bm a, \bm b]} f(\bm x),\quad f(\bm d) = \inf_{\bm x \in [\bm a, \bm b]} f(\bm x). $$
        \end{itemize}
    \end{theoremenv}
    \begin{definitionenv}
        Suppose $\left[ \bm a, \bm b \right]$ and $P_k$ is a partition of $[a_k, b_k]$ for $k=1, \ldots, n$. We call $P := P_1 \times \cdots \times P_n$ a \textbf{partition} of $[\bm a, \bm b]$. 
    \end{definitionenv}
    \begin{definitionenv}
        Let $f$ be a continuous scalar field on $[\bm a, \bm b] \subseteq \R^n$ and let $M(f)$ and $m(f)$ denote the absolute maximum and minimum values of $f$ on $[\bm a, \bm b]$, respectively. We call $M(f) - m(f)$ the \textbf{span} of $f$ on $[\bm a, \bm b]$.
    \end{definitionenv}
    \begin{theoremenv}\label{thm:small-span}
        Suppose $f$ is a scalar field continuous on $[\bm a, \bm b]$. Then, for every $\varepsilon > 0$, there is a partition of $[\bm a, \bm b]$ such that the span of $f$ on every subinterval of this partition is less than $\varepsilon$.
    \end{theoremenv}

    \newpage